\lesson{3}{The gravity droop}{A proportional controller can only push when there is an error — so it keeps one.}{1}{droop}{Part I · PID}
\label{les:droop}

\lead{Take the feedforward away. The controller no longer knows that the drone has weight; it only has its error and its gains, and the integral is switched off. The drone takes off, climbs — and stops a metre short of the setpoint, hovering there perfectly calm, perfectly wrong. This lesson finds out exactly how wrong, and why, and then fixes it in two different ways.}

\section{The loop without feedforward}

In this lesson the thrust is made of the proportional and derivative terms only:
\begin{equation}
  T = K_p\,e - K_d\,\dot y ,\qquad e = r - y .
\end{equation}
(The D term here acts on the measured velocity rather than on $\dot e$; for a constant setpoint the two are the same, $\dot e = -\dot y$, and we will see in \lessonref{les:kick} why this form is preferred.) Substituting into the plant \eqref{eq:meet-plant} with no wind:
\begin{equation}\label{eq:droop-loop}
  m\ddot y = K_p\,(r - y) - K_d\,\dot y - m g .
\end{equation}
The weight no longer cancels. Something must carry it.

\section{Where can the drone rest?}

We are not asking yet how the drone moves, only where it can stop. Such a point has its own name.

\begin{definition}{Equilibrium and steady state}
An \term{equilibrium} of a dynamical system is a state in which it can remain forever: all time derivatives are zero. If the system is stable, it approaches an equilibrium after a disturbance or a setpoint change; the values it approaches are its \term{steady state}, written with the subscript \emph{ss}.
\end{definition}

\begin{example}[The droop, step by step]
\label{ex:droop-ess}
Find the equilibrium of \cref{eq:droop-loop} for $m = \SI{1}{kg}$, $K_p = 10$, $r = \SI{2}{m}$.

\textbf{Step 1 — set the derivatives to zero.} At rest $\dot y = 0$ and $\ddot y = 0$. The left-hand side vanishes, and so does the D term:
$0 = K_p(r - y_{ss}) - mg$.

\textbf{Step 2 — solve for the error.} $K_p\,e_{ss} = mg$, so
\begin{equation}\label{eq:droop-ess}
  e_{ss} = \frac{mg}{K_p} = \frac{1 \cdot 9.81}{10} = \SI{0.981}{m}.
\end{equation}

\textbf{Step 3 — the altitude.} $y_{ss} = r - e_{ss} = 2 - 0.981 = \SI{1.019}{m}$.

\textbf{Step 4 — the forces at rest.} The P term is $K_p e_{ss} = 10 \cdot 0.981 = \SI{9.81}{N}$ — exactly the weight. The D term is zero. The P term carries the whole drone.

\textbf{Step 5 — a stiffer spring.} With $K_p = 40$: $e_{ss} = 9.81/40 = \SI{0.245}{m}$, $y_{ss} = \SI{1.755}{m}$. Four times the gain, a quarter of the droop, but still a droop.
\end{example}

\Cref{fig:droop-run} confirms the prediction to the centimetre: the simulator's drone settles at \SI{1.019}{m} with $K_p = 10$ and at \SI{1.755}{m} with $K_p = 40$. On the right, the P term ends at \SI{9.81}{N} in both runs: it is doing the feedforward's job.

\simfig{droop_run}{Left: the altitude with P (and D) only, no feedforward, for two stiffnesses; the dotted lines are the predictions of Example~\ref{ex:droop-ess}. Right: the P term (dashed: $K_p=10$, solid: $K_p=40$) ends at \SI{9.81}{N} in both cases — the whole weight hangs on the spring.}{fig:droop-run}

\marginscreen{1905 675 0 112}{droop}{The lesson panel in the app, with the same prediction.}{fig:droop-panel}

There is a simple way to see why the droop is unavoidable. The drone is in a tug of war between gravity, which never tires, and the P term, which pulls only as hard as the error is large. The rope settles where the two pulls are equal — and since gravity pulls with \SI{9.81}{N}, the error must be large enough for P to pull with \SI{9.81}{N} too. Zero error would mean zero pull, and the drone would fall.

\begin{keyidea}[Steady-state error of a proportional loop]
Under a constant load $F$ that nothing else compensates, a proportional loop settles with the error
\begin{equation*}
  e_{ss} = F / K_p .
\end{equation*}
Stiffer springs sag less, but no finite stiffness removes the sag. The same holds for a constant wind (the load is the drag) or a wrong mass estimate (the load is the part of the weight that the feedforward misses).
\end{keyidea}

\subsection{Why not just raise the gain?}
Equation~\eqref{eq:droop-ess} suggests an easy way out: make $K_p$ huge. To keep the droop below \SI{1}{cm} we would need $K_p = 9.81/0.01 = \SI{981}{N/m}$. In \lessonref{les:spring} we saw the first problem with large gains — the motors saturate and the loop oscillates harder. \Lessonref{les:edge} will show the fundamental one: the motors lag and the controller samples, and together they put an upper limit on any gain, around $K_p \approx 200$ for this drone. There are two better ways.

\section{Cure one: remember the error}

The integral term accumulates the error:
\begin{equation}
  I(t) = K_i\int_0^t e(\tau)\,\dd\tau ,\qquad \dot I = K_i\,e .
\end{equation}
As long as the drone hangs low, $e > 0$ and $I$ grows. It keeps growing until it can carry the weight by itself; only then can the error — and P — return to zero.

\begin{result}[The integral can only rest at zero error]
At any equilibrium all derivatives vanish, in particular $\dot I = K_i\,e_{ss} = 0$. With $K_i \neq 0$ this forces $e_{ss} = 0$. The integral then holds the whole load: $I_{ss} = mg$ here.
\end{result}

\Lessonref{les:integral} is devoted to this idea; here we watch it work (\cref{fig:droop-ion}) and ask how long it takes.

\simfig{droop_ion}{At $t = \SI{10}{s}$ the integral is switched on. It slowly takes over the weight from the P term (bottom); the altitude creeps up to the setpoint (top). With the default $K_i = 0.8$ this takes a long time.}{fig:droop-ion}

\subsection{How fast does the integral pay back the load?}
The creep in \cref{fig:droop-ion} is slow: the drone moves a few centimetres per second at most. That suggests a simplification: treat it as \emph{quasi-static}, i.e.\ at every moment the drone is almost at rest, so its acceleration and velocity are negligible and the forces almost balance.

\begin{example}[The integral time constant]
\label{ex:droop-tau}
Derive how the error decays after the integral is switched on, and check the prediction against the simulator ten seconds later.

\textbf{Step 1 — the quasi-static balance.} With $\ddot y \approx 0$ and $\dot y \approx 0$, the forces balance: $K_p\,e + I = mg$.

\textbf{Step 2 — differentiate.} The weight is constant, so $K_p\,\dot e + \dot I = 0$.

\textbf{Step 3 — insert the integral's law.} $\dot I = K_i e$, hence $K_p\,\dot e = -K_i\,e$, or
\begin{equation*}
  \dot e = -\frac{K_i}{K_p}\,e .
\end{equation*}

\textbf{Step 4 — solve.} This is the simplest differential equation there is: the rate of change is proportional to the value. Separating variables, $\dd e/e = -(K_i/K_p)\,\dd t$, and integrating from the moment $t_0$ when the integral is switched on,
\begin{equation}\label{eq:droop-tau}
  e(t) = e_{ss}\,e^{-(t-t_0)/\tau_I},\qquad \tau_I = \frac{K_p}{K_i}.
\end{equation}

\textbf{Step 5 — numbers.} $\tau_I = 10/0.8 = \SI{12.5}{s}$. Ten seconds after switching on: $e = 0.981\,e^{-10/12.5} = 0.981 \cdot 0.449 = \SI{0.44}{m}$. The simulator shows \SI{0.447}{m}.

\textbf{Step 6 — when is it done?} To shrink below \SI{1}{cm} the error must fall by a factor of 98: $e^{-t/12.5} = 1/98$, $t = 12.5\ln 98 = \SI{57}{s}$. With $K_i = 3$: $\tau_I = \SI{3.3}{s}$, done in \SI{15}{s}.
\end{example}

The time constant $K_p/K_i$ is exactly the \term{integral time} $T_i$ of the \enquote{standard form} of the PID that industrial controllers use,\aside{Standard form: $u = K_p\big(e + \tfrac{1}{T_i}\int e + T_d\,\dot e\big)$ with $T_i = K_p/K_i$ and $T_d = K_d/K_p$. It is the same controller with different knobs.} and it tells you how long the integral needs to pay back a load. It also explains why the quasi-static assumption worked: \SI{12.5}{s} is much slower than the spring–damper motion of the drone, which settles within a second or two.

\section{Cure two: know the load}

The other cure is not to wait for an error at all. The weight is known — roughly. So add it to the command directly:
\begin{equation}\label{eq:droop-ff}
  T = \underbrace{\hat m\,g}_{\text{feedforward}} + \underbrace{K_p\,e + I - K_d\,\dot y}_{\text{feedback}} .
\end{equation}

\begin{definition}{Feedforward}
A \term{feedforward} term is a part of the command computed from a model of the plant and from known inputs (the setpoint, a measured load), \emph{not} from the error. It acts immediately, without waiting for an error to appear, and it is exactly as good as the model.
\end{definition}

Now the feedback only has to correct the \emph{difference} between the real weight and the assumed one.

\begin{example}[Feedforward with a wrong mass]
\label{ex:droop-ff}
The feedforward uses $\hat m = \SI{0.9}{kg}$ for the \SI{1}{kg} drone; I is off, $K_p = 10$. Where does the drone settle?

\textbf{Step 1 — the equilibrium.} With \cref{eq:droop-ff} at rest: $0 = \hat m g + K_p e_{ss} - m g$.

\textbf{Step 2 — the error.} $e_{ss} = (m - \hat m)g/K_p = 0.1 \cdot 9.81/10 = \SI{0.098}{m}$.

\textbf{Step 3 — compare.} Without feedforward the droop was \SI{0.98}{m}; with a feedforward that is 10\,\% wrong it is \SI{0.098}{m}, ten times smaller. The feedforward removes the part of the load it knows; feedback deals with the rest.
\end{example}

\begin{keyidea}[Feedforward and feedback]
Feedforward acts on what you know, before any error exists. Feedback acts on what you do not know, after it has caused an error. Good controllers use both: the feedforward carries the bulk of the known load, the feedback (and especially its integral) mops up the model errors. Every flight controller in this book, up to the trajectory planners of Part~II, is built this way.
\end{keyidea}

\begin{inapp}
\begin{itemize}[leftmargin=1.2em]
  \item The lesson sets \param{control.feedforward = false} and \param{control.alt.iOn = false}. The switch is \ui{Controller\then Gravity feedforward}; the assumed mass $\hat m$ is \param{control.model.mass}.
  \item The feedforward is added in \src{src/control/altitude.ts}: the altitude loop passes $\hat m g$ to the PID as its \texttt{ff} argument, and the PID reports it as its own coloured part, \FFt.
  \item The info card always prints the \ui{predicted e\_ss}: the part of the weight the feedforward does not carry, divided by $K_p$ — $mg/K_p$ without feedforward, $(m - \hat m)g/K_p$ with it — or \enquote{0 (I term on)}. It is computed in \src{src/ui/charts/InfoCard.tsx}, exactly as in Examples~\ref{ex:droop-ess} and \ref{ex:droop-ff}.
\end{itemize}
\end{inapp}

\begin{trythis}
\begin{enumerate}
  \item Raise $K_p$ to 40: the droop shrinks to about \SI{0.25}{m} but never vanishes.
  \item Back to 10. Switch \It on and time the creep. Compare with Example~\ref{ex:droop-tau}. Then try $K_i = 3$.
  \item Switch \It off again and the feedforward on instead: the droop disappears at once.
\end{enumerate}
\predict{with the feedforward on but $\hat m = \SI{0.9}{kg}$ (and I off), where does the drone settle? (Example~\ref{ex:droop-ff}.)}
\end{trythis}

\begin{curiosity}{The grid that sags on purpose}
\photo{turbinehall}{A turbine hall of the 1920s. Every generator in a power grid runs a speed governor — with a deliberate droop.}
The electrical grid of a continent is one enormous machine turning in synchrony: every generator in Europe spins in step with a frequency of \SI{50}{Hz}. When you switch on a kettle, the load rises, the turbines slow down a little, and the frequency dips. Each power station answers with a speed governor that opens its steam or water valve.

Here is the surprise: those governors are \emph{proportional} controllers, and their droop is deliberate. A typical setting is 4–5\,\%: a generator increases its output by its full rated power if the frequency falls by 4–5\,\%. Why not an integral, which would restore \SI{50}{Hz} exactly? Because hundreds of generators share the load. If each of them integrated its own error, they would fight over it forever; with proportional droop every machine takes a share proportional to its size, and the sharing is stable.

The small remaining frequency error is then removed on a slower time scale, over tens of seconds to minutes, by a single central \emph{integral} loop, the secondary control, which asks selected plants to raise their setpoints. It is the structure of this lesson, spread over a continent: a fast proportional loop that accepts a droop, and a slow integral that pays it back.
\end{curiosity}

\begin{remember}
\begin{itemize}
  \item An equilibrium is found by setting all derivatives to zero.
  \item A proportional loop under a constant load $F$ settles with $e_{ss} = F/K_p$. Raising $K_p$ reduces the droop but cannot remove it.
  \item An integral can only rest at zero error, so it removes any constant droop — slowly, with time constant about $K_p/K_i$.
  \item Feedforward removes the known part of the load immediately; feedback handles what the model gets wrong.
\end{itemize}
\end{remember}

\begin{exercises}
\exercise[1] The drone carries a \SI{200}{g} camera that the controller does not know about ($\hat m = \SI{1}{kg}$, feedforward on, I off, $K_p = 10$). Where does it hover when asked for \SI{2}{m}?
\answer{The missing force is $0.2 \cdot 9.81 = \SI{1.96}{N}$, so $e_{ss} = \SI{0.196}{m}$: it hovers at about \SI{1.80}{m}.}
\exercise[1] A steady wind blows horizontally on the L2 point-mass drone with a drag force of \SI{3}{N}. With a P-only horizontal controller of $K_p = 4$, how far downwind does it settle? Which $K_p$ would be needed for \SI{10}{cm}?
\answer{$e_{ss} = 3/4 = \SI{0.75}{m}$; $K_p = 30$.}
\exercise[1] With $K_p = 10$ and $K_i = 0.8$, how long after switching the integral on is the droop halved?
\answer{$t = \tau_I\ln2 = 12.5\cdot0.693 = \SI{8.7}{s}$.}
\exercise[2] Derive \cref{eq:droop-tau} again, this time keeping the D term and the inertia: write the closed-loop characteristic polynomial $m s^3 + K_d s^2 + K_p s + K_i$ and show that for small $K_i$ it has one real root close to $-K_i/K_p$. (Hint: substitute $s = -\varepsilon$ and keep first-order terms.)
\answer{$-m\varepsilon^3 + K_d\varepsilon^2 - K_p\varepsilon + K_i \approx -K_p\varepsilon + K_i = 0$, so $\varepsilon \approx K_i/K_p$.}
\exercise[2] In the power-grid story, two generators of \SI{100}{MW} and \SI{300}{MW} both have a 5\,\% droop. A load of \SI{40}{MW} is added. How does it split, and how far does the frequency fall (before secondary control acts)?
\answer{\SI{10}{MW} and \SI{30}{MW}; $\Delta f/f_0 = 0.05\cdot40/400 = 0.5\,\%$, i.e.\ \SI{0.25}{Hz}.}
\exercise[3] The quasi-static assumption of Example~\ref{ex:droop-tau} neglects $m\ddot y$ and $K_d\dot y$. Using the solution \eqref{eq:droop-tau}, estimate both at the moment the integral is switched on and compare them with the \SI{9.81}{N} being shared. Is the assumption justified?
\answer{$\dot y = e_{ss}/\tau_I = \SI{0.078}{m/s}$, $K_d\dot y = \SI{0.55}{N}$, about 6\,\%; $m\ddot y \approx \SI{0.006}{N}$. Reasonable.}
\end{exercises}
